\chapter{2019年天津卷（理）}

\section{单选题}

\begin{problem}
设集合 \( A=\{-1,1,2,3,5\} \)，\( B=\{2,3,4\} \)，\( C=\{x\in \R\mid 1\le x<3\} \)，则 \( (A\cap C)\cup B= \)（\quad）

\choices
    {\(\{2\}\)}
    {\(\{2,3\}\)}
    {\(\{-1,2,3\}\)}
    {\(\{1,2,3,4\}\)}

\end{problem}

\begin{answer}
D
\end{answer}

\begin{solution}
因为 \( A\cap C=\{1,2\} \)，

所以
\((A\cap C)\cup B=\{1,2,3,4\}\).

故选 \( D \).
\end{solution}
\begin{problem}
设变量 \( x,y \) 满足约束条件 \( \left\{\begin{array}{l}
x+y-2\le 0,\\
x-y+2\ge 0,\\
x\ge -1,\\
y\ge -1,
\end{array}\right. \)
则目标函数 \( z=-4x+y \) 的最大值为（\quad）

\choices
    {\(2\)}
    {\(3\)}
    {\(5\)}
    {\(6\)}

\end{problem}

\begin{answer}
C
\end{answer}

\begin{solution}
已知不等式组表示的平面区域如图中的阴影部分.

\begin{center}
\bitmapinclude[max width=5.2cm,max totalheight=4.8cm]{img/2019/tianjin_science/q2_fig1.png}
\end{center}

目标函数的几何意义是直线 \( y=4x+z \) 在 \( y \) 轴上的截距，

故目标函数在点 \( A \) 处取得最大值.

由 \( x-y+2=0 \) 与 \( x=-1 \) 得 \( A(-1,1) \)，

所以
\(z_{\max}=-4\times(-1)+1=5\).

故选 \( C \).
\end{solution}
\begin{problem}
设 \( x\in \R \)，则“\( x^2-5x<0 \)”是“\( |x-1|<1 \)”的（\quad）

\choices
    {充分而不必要条件}
    {必要而不充分条件}
    {充要条件}
    {既不充分也不必要条件}

\end{problem}

\begin{answer}
B
\end{answer}

\begin{solution}
由 \( x^2-5x<0 \) 得
\(0<x<5\).

由 \( |x-1|<1 \) 得
\(0<x<2\).

故 \( 0<x<5 \) 推不出 \( |x-1|<1 \)，而 \( |x-1|<1 \) 能推出 \( 0<x<5 \).

所以“\( x^2-5x<0 \)”是“\( |x-1|<1 \)”的必要而不充分条件.

故选 \( B \).
\end{solution}
\begin{problem}
阅读下边的程序框图，运行相应的程序，输出 \( S \) 的值为（\quad）

\begin{center}
\bitmapinclude[max width=4.2cm,max totalheight=4.8cm]{img/2019/tianjin_science/q4_fig1.png}
\end{center}

\choices
    {\(5\)}
    {\(8\)}
    {\(24\)}
    {\(29\)}

\end{problem}

\begin{answer}
B
\end{answer}

\begin{solution}
根据程序框图，逐步写出运算结果：
\(S=1\)，\(i=2\to j=1\)，\(S=1+2\cdot 2^1=5\)，\(i=3\)，\(S=8\)，\(i=4\).

结束循环，故输出 \( 8 \).

故选 \( B \).
\end{solution}
\begin{problem}
已知抛物线 \( y^2=4x \) 的焦点为 \( F \)，准线为 \( l \)，若 \( l \) 与双曲线 \( \dfrac{x^2}{a^2}-\dfrac{y^2}{b^2}=1 \)（\(a>0\)，\(b>0\)）的两条渐近线分别交于点 \( A \) 和点 \( B \)，且 \( |AB|=4|OF| \)（\( O \) 为原点），则双曲线的离心率为（\quad）

\choices
    {\(\sqrt2\)}
    {\(\sqrt3\)}
    {\(2\)}
    {\(\sqrt5\)}

\end{problem}

\begin{answer}
D
\end{answer}

\begin{solution}
抛物线 \( y^2=4x \) 的准线方程为 \( x=-1 \).

双曲线的渐近线方程为
\( y=\pm \frac{b}{a}x \).
故
\(A\left( -1,\frac{b}{a} \right)\)，\(B\left( -1,-\frac{b}{a} \right)\).

所以
\( |AB|=\frac{2b}{a} \).

又 \( |OF|=1 \)，且 \( |AB|=4|OF| \)，所以
\( \frac{2b}{a}=4 \)，
即
\( b=2a \).

因此双曲线的离心率
\( e=\frac{c}{a}=\frac{\sqrt{a^2+b^2}}{a}=\sqrt5 \).

故选 \( D \).
\end{solution}
\begin{problem}
已知 \( a=\log_5 2 \)，\( b=\log_{0.5}0.2 \)，\( c=0.5^{0.2} \)，则 \( a,b,c \) 的大小关系为（\quad）

\choices
    {\(a<c<b\)}
    {\(a<b<c\)}
    {\(b<c<a\)}
    {\(c<a<b\)}

\end{problem}

\begin{answer}
A
\end{answer}

\begin{solution}
因为
\(a=\log_5 2<\log_5 \sqrt5=\frac12\)，

又
\(b=\log_{0.5}0.2>\log_{0.5}0.25=2\)，

且
\(0.5^1<0.5^{0.2}<0.5^0\)，
所以
\(\frac12<c<1\).

因此
\(a<c<b\).

故选 \( A \).
\end{solution}
\begin{problem}
已知函数 \( f(x)=A\sin(\omega x+\varphi) \)（\(A>0\)，\(\omega>0\)，\(|\varphi|<\pi\)）是奇函数，将 \( y=f(x) \) 的图象上所有点的横坐标伸长到原来的 \( 2 \) 倍（纵坐标不变），所得图象对应的函数为 \( g(x) \). 若 \( g(x) \) 的最小正周期为 \( 2\pi \)，且 \( g\!\left( \dfrac{\pi}{4} \right)=\sqrt2 \)，则 \( f\!\left( \dfrac{3\pi}{8} \right)= \)（\quad）

\choices
    {\(-2\)}
    {\(-\sqrt2\)}
    {\(\sqrt2\)}
    {\(2\)}

\end{problem}

\begin{answer}
C
\end{answer}

\begin{solution}
因为 \( f(x) \) 为奇函数，所以
\(f(0)=A\sin\varphi=0\).

又 \( |\varphi|<\pi \)，所以
\(\varphi=0\).

把其图象上各点的横坐标伸长到原来的 \( 2 \) 倍，得
\(g(x)=A\sin \frac{\omega x}{2}\).

由 \( g(x) \) 的最小正周期为 \( 2\pi \) 可得
\(\omega=2\).

由
\(g\!\left( \frac{\pi}{4} \right)=\sqrt2\)
可得
\(A=2\).

所以
\(f(x)=2\sin 2x\)，
从而
\(f\!\left( \frac{3\pi}{8} \right)=2\sin \frac{3\pi}{4}=\sqrt2\).

故选 \( C \).
\end{solution}
\begin{problem}
已知 \( a\in \R \)，设函数 \( f(x)=\left\{\begin{array}{ll}
x^2-2ax+2a, & x\le 1,\\
x-a\ln x, & x>1.
\end{array}\right. \)
若关于 \( x \) 的不等式 \( f(x)\ge 0 \) 在 \( \R \) 上恒成立，则 \( a \) 的取值范围为（\quad）

\choices
    {\([0,1]\)}
    {\([0,2]\)}
    {\([0,\e]\)}
    {\([1,\e]\)}

\end{problem}

\begin{answer}
C
\end{answer}

\begin{solution}
首先 \( f(0)\ge 0 \)，即
\(a\ge 0\).

当 \( 0\le a\le 1 \) 时，
\(f(x)=x^2-2ax+2a=(x-a)^2+2a-a^2\ge 2a-a^2=a(2-a)>0\).

当 \( a<1 \) 时，
\(f(1)=1>0\).

故当 \( a\ge 0 \) 时，\( x^2-2ax+2a\ge 0 \) 在 \( (-\infty,1] \) 上恒成立.

若 \( x-a\ln x\ge 0 \) 在 \( (1,+\infty) \) 上恒成立，则
\(a\le \frac{x}{\ln x}\)
在 \( (1,+\infty) \) 上恒成立.

令
\(g(x)=\frac{x}{\ln x}\)，
则
\(g'(x)=\frac{\ln x-1}{(\ln x)^2}\).

因为当 \( 1<x<\e \) 时，\( \ln x-1<0 \)，所以 \( g'(x)<0 \)；当 \( x>\e \) 时，\( \ln x-1>0 \)，所以 \( g'(x)>0 \). 因此 \( x=\e \) 为函数 \( g(x) \) 在 \( (1,+\infty) \) 上唯一的极小值点，也是最小值点，

故
\(g(x)_{\min}=g(\e)=\e\)，
所以
\(a\le \e\).

综上可知，\( a \) 的取值范围是 \( [0,\e] \).

故选 \( C \).
\end{solution}
\section{填空题}

\begin{problem}
\( i \) 是虚数单位，则 \( \left| \dfrac{5-i}{1+i} \right|=\fillinblank{} \).

\end{problem}

\begin{answer}
\(\sqrt{13}\)
\end{answer}

\begin{solution}
解法一：
\(\left| \frac{5-i}{1+i} \right|
=\left| \frac{(5-i)(1-i)}{(1+i)(1-i)} \right|
=|2-3i|
=\sqrt{13}\).

解法二：
\(\left| \frac{5-i}{1+i} \right|
=\frac{|5-i|}{|1+i|}
=\frac{\sqrt{26}}{\sqrt2}
=\sqrt{13} \).
\end{solution}
\begin{problem}
\( \left( 2x-\dfrac{1}{8x^3} \right)^8 \) 的展开式中的常数项为 \fillinblank{}.

\end{problem}

\begin{answer}
28
\end{answer}

\begin{solution}
展开式的通项为
\(T_{r+1}=C_8^r(2x)^{8-r}\left( -\frac{1}{8x^3} \right)^r
=(-1)^r2^{8-4r}C_8^r x^{8-4r} \).

由
\(8-4r=0\)
得
\(r=2\).

故所求的常数项为
\((-1)^2C_8^2=28\).
\end{solution}
\begin{problem}
已知四棱锥的底面是边长为 \( \sqrt2 \) 的正方形，侧棱长均为 \( \sqrt5 \). 若圆柱的一个底面的圆周经过四棱锥四条侧棱的中点，另一个底面的圆心为四棱锥底面的中心，则该圆柱的体积为 \fillinblank{}.

\end{problem}

\begin{answer}
\(\dfrac{\pi}{4}\)
\end{answer}

\begin{solution}
四棱锥的高为
\(\sqrt{5-1}=2\).

故圆柱的高为 \( 1 \)，圆柱的底面半径为
\( \frac12 \).

故其体积为
\( V=\pi \times \left( \frac12 \right)^2 \times 1=\frac{\pi}{4} \).
\end{solution}
\begin{problem}
设 \( a\in \R \)，直线 \( ax-y+2=0 \) 和圆 \( \left\{\begin{array}{l}
x=2+2\cos\theta,\\
y=1+2\sin\theta
\end{array}\right. \)（\( \theta \) 为参数）相切，则 \( a \) 的值为 \fillinblank{}.

\end{problem}

\begin{answer}
\(\dfrac{3}{4}\)
\end{answer}

\begin{solution}
由圆的参数方程可知圆心坐标为 \( (2,1) \)，半径为 \( 2 \).

因为直线与圆相切，所以圆心到直线的距离等于半径，即
\(\frac{|2a+1|}{\sqrt{a^2+1}}=2\).

化简得
\(4a^2+4a+1=4a^2+4\).

解得
\(a=\frac34\).
\end{solution}
\begin{problem}
设 \( x>0 \)，\( y>0 \)，\( x+2y=5 \)，则 \( \dfrac{(x+1)(2y+1)}{\sqrt{xy}} \) 的最小值为 \fillinblank{}.

\end{problem}

\begin{answer}
\(4\sqrt3\)
\end{answer}

\begin{solution}
因为
\(\frac{(x+1)(2y+1)}{\sqrt{xy}}
=\frac{2xy+x+2y+1}{\sqrt{xy}}
=\frac{2xy+6}{\sqrt{xy}}
\ge \frac{2\sqrt{2xy\cdot 6}}{\sqrt{xy}}
=4\sqrt3\).

当且仅当 \( xy=3 \)，即 \(x=3\)，\(y=1\) 时等号成立.

故所求的最小值为 \( 4\sqrt3 \).
\end{solution}
\begin{problem}
在四边形 \( ABCD \) 中，\( AD//BC \)，\( AB=2\sqrt3 \)，\( AD=5 \)，\( \angle A=30^\circ \)，点 \( E \) 在线段 \( CB \) 的延长线上，且 \( AE=BE \)，则 \( \overrightarrow{BD}\cdot \overrightarrow{AE}=\fillinblank{} \).

\begin{center}
\bitmapinclude[max width=7.8cm,max totalheight=4.8cm]{img/2019/tianjin_science/q14_fig1.png}
\end{center}

\end{problem}

\begin{answer}
\(-1\)
\end{answer}

\begin{solution}
解法一：如图，过点 \( B \) 作 \( AE \) 的平行线交 \( AD \) 于 \( F \).

因为 \( AE=BE \)，故四边形 \( AEBF \) 为菱形.

因为 \( \angle BAD=30^\circ \)，\( AB=2\sqrt3 \)，所以 \( AF=2 \)，即
\(\overrightarrow{AF}=\frac25\overrightarrow{AD}\).

因为
\(\overrightarrow{AE}=\overrightarrow{FB}=\overrightarrow{AB}-\overrightarrow{AF}
=\overrightarrow{AB}-\frac25\overrightarrow{AD}\)，
所以
\(\overrightarrow{BD}\cdot \overrightarrow{AE}
=(\overrightarrow{AD}-\overrightarrow{AB})\cdot \left( \overrightarrow{AB}-\frac25\overrightarrow{AD} \right)
=\frac75\overrightarrow{AB}\cdot \overrightarrow{AD}-AB^2-\frac25AD^2
=\frac75\times 2\sqrt3\times 5\times \frac{\sqrt3}{2}-12-10=-1\).

解法二：建立如图所示的直角坐标系.

\bitmapfigure[max width=7.8cm,max totalheight=4.8cm]{img/2019/tianjin_science/q14_solution_fig1.png}

则 \(B(2\sqrt3,0)\)，\(D\left( \frac{5\sqrt3}{2},\frac52 \right)\).

因为 \( AD//BC \)，\( \angle BAD=30^\circ \)，所以 \( \angle CBE=30^\circ \).

因为 \( AE=BE \)，所以 \( \angle BAE=30^\circ \).

所以直线 \( BE \) 的斜率为 \( \dfrac{\sqrt3}{3} \)，其方程为
\(y=\frac{\sqrt3}{3}(x-2\sqrt3)\)，
直线 \( AE \) 的斜率为 \( -\dfrac{\sqrt3}{3} \)，其方程为
\(y=-\frac{\sqrt3}{3}x\).

由 \(\frac{\sqrt3}{3}(x-2\sqrt3)=-\frac{\sqrt3}{3}x\)
得
\(x=\sqrt3\)，\(y=-1\)，
所以
\(E(\sqrt3,-1)\).

因此
\(\overrightarrow{BD}=\left( \frac{\sqrt3}{2},\frac52 \right)\)，\(\overrightarrow{AE}=(\sqrt3,-1)\)，
从而
\(\overrightarrow{BD}\cdot \overrightarrow{AE}=-1\).
\end{solution}
\section{解答题}

\begin{problem}
在 \( \triangle ABC \) 中，内角 \( A,B,C \) 所对的边分别为 \( a,b,c \). 已知 \( b+c=2a \)，\( 3c\sin B=4a\sin C \).

\begin{enumerate}
\item 求 \( \cos B \) 的值；
\item 求 \( \sin \left( 2B+\dfrac{\pi}{6} \right) \) 的值.
\end{enumerate}

\end{problem}

\begin{solution}
（1）在 \( \triangle ABC \) 中，由正弦定理
\(\frac{b}{\sin B}=\frac{c}{\sin C}\)，
得
\(b\sin C=c\sin B\).

又由 \( 3c\sin B=4a\sin C \)，得
\(3b\sin C=4a\sin C\)，
即
\(3b=4a\).

又因为 \( b+c=2a \)，所以
\(b=\frac43a\)，\(c=\frac23a\).

由余弦定理，
\(\cos B=\frac{a^2+c^2-b^2}{2ac}
=\frac{a^2+\frac49a^2-\frac{16}{9}a^2}{2\cdot a\cdot \frac23a}
=-\frac14\).

（2）由（1）得
\(\sin B=\sqrt{1-\cos^2 B}=\frac{\sqrt{15}}{4}\).

从而
\(\sin 2B=2\sin B\cos B=-\frac{\sqrt{15}}{8}\)，
\(\cos 2B=\cos^2 B-\sin^2 B=-\frac78\).

所以
\(\sin \left( 2B+\frac{\pi}{6} \right)
=\sin 2B\cos \frac{\pi}{6}+\cos 2B\sin \frac{\pi}{6}
=-\frac{\sqrt{15}}{8}\cdot \frac{\sqrt3}{2}-\frac78\cdot \frac12
=-\frac{3\sqrt5+7}{16}\).
\end{solution}
\begin{problem}
设甲、乙两位同学上学期间，每天 \( 7{:}30 \) 之前到校的概率均为 \( \dfrac{2}{3} \). 假定甲、乙两位同学到校情况互不影响，且任一同学每天到校情况相互独立.

\begin{enumerate}
\item 用 \( X \) 表示甲同学上学期间的三天中 \( 7{:}30 \) 之前到校的天数，求随机变量 \( X \) 的分布列和数学期望；
\item 设 \( M \) 为事件“上学期间的三天中，甲同学在 \( 7{:}30 \) 之前到校的天数比乙同学在 \( 7{:}30 \) 之前到校的天数恰好多 \( 2 \)”，求事件 \( M \) 发生的概率.
\end{enumerate}

\end{problem}

\begin{solution}
（1）因为甲同学上学期间的三天中到校情况相互独立，且每天 \( 7{:}30 \) 之前到校的概率均为 \( \dfrac{2}{3} \)，

故
\(X\sim B\left( 3,\frac23 \right)\)，
从而
\(P(X=k)=C_3^k\left( \frac23 \right)^k \left( \frac13 \right)^{3-k}\)（\(k=0,1,2,3\)）.

所以，随机变量 \( X \) 的分布列为
\examdisplayarray{}{c|cccc}{
X & 0 & 1 & 2 & 3 \\
\hline
P & \dfrac{1}{27} & \dfrac{2}{9} & \dfrac{4}{9} & \dfrac{8}{27}
}{}

随机变量 \( X \) 的数学期望为
\(E(X)=3\times \frac23=2\).

（2）设乙同学上学期间的三天中 \( 7{:}30 \) 之前到校的天数为 \( Y \)，则
\(Y\sim B\left( 3,\frac23 \right)\).

且
\(M=\{X=3,Y=1\}\cup \{X=2,Y=0\}\).

由题意知事件 \( \{X=3,Y=1\} \) 与 \( \{X=2,Y=0\} \) 互斥，

且事件 \( \{X=3\} \) 与 \( \{Y=1\} \)、事件 \( \{X=2\} \) 与 \( \{Y=0\} \) 均相互独立，

从而
\(P(M)=P(\{X=3,Y=1\}\cup \{X=2,Y=0\})
=P(X=3,Y=1)+P(X=2,Y=0)
=P(X=3)P(Y=1)+P(X=2)P(Y=0)
=\frac{8}{27}\times \frac29+\frac49\times \frac{1}{27}
=\frac{20}{243}\).
\end{solution}
\begin{problem}
如图，\( AE\perp \)平面 \( ABCD \)，\( CF//AE \)，\( AD//BC \)，\( AD\perp AB \)，\( AB=AD=1 \)，\( AE=BC=2 \).

\begin{enumerate}
\item 求证：\( BF// \)平面 \( ADE \)；
\item 求直线 \( CE \) 与平面 \( BDE \) 所成角的正弦值；
\item 若二面角 \( E-BD-F \) 的余弦值为 \( \dfrac{1}{3} \)，求线段 \( CF \) 的长.
\end{enumerate}

\begin{center}
\bitmapinclude[max width=8.5cm,max totalheight=4.8cm]{img/2019/tianjin_science/q17_fig1.png}
\end{center}
\end{problem}

\begin{solution}
（1）依题意，可以建立以 \( A \) 为原点，分别以 \( AB \)、\( AD \)、\( AE \) 的方向为 \( x \) 轴、\( y \) 轴、\( z \) 轴正方向的空间直角坐标系.

\bitmapfigure[max width=7.8cm,max totalheight=4.8cm]{img/2019/tianjin_science/q17_solution_fig2.png}

可得
\(A(0,0,0)\)，\(B(1,0,0)\)，\(C(1,2,0)\)，\(D(0,1,0)\)，\(E(0,0,2)\).

设 \( CF=h \)（\( h>0 \)），则
\(F(1,2,h)\).

依题意，\(\overrightarrow{AB}=(1,0,0)\) 是平面 \( ADE \) 的法向量，

又
\(\overrightarrow{BF}=(0,2,h)\)，
可得
\(\overrightarrow{BF}\cdot \overrightarrow{AB}=0\).

又因为直线 \( BF\not\subset \)平面 \( ADE \)，所以
\(BF// \text{平面 } ADE\).

（2）依题意，
\(\overrightarrow{BD}=(-1,1,0)\)，\(\overrightarrow{BE}=(-1,0,2)\)，\(\overrightarrow{CE}=(-1,-2,2)\).

设 \( \bs{n}=(x,y,z) \) 为平面 \( BDE \) 的法向量，则
\(\bs{n}\cdot \overrightarrow{BD}=0\) 且 \(\bs{n}\cdot \overrightarrow{BE}=0\)，
即 \(-x+y=0\)，\(-x+2z=0\).

不妨令 \( z=1 \)，可得
\(\bs{n}=(2,2,1)\).

因此
\(\cos \angle \left( \overrightarrow{CE},\bs{n} \right)=\frac{\overrightarrow{CE}\cdot \bs{n}}{|\overrightarrow{CE}||\bs{n}|}=-\frac49\).

所以，直线 \( CE \) 与平面 \( BDE \) 所成角的正弦值为
\(\frac49\).

（3）设 \( \bs{m}=(x,y,z) \) 为平面 \( BDF \) 的法向量，则
\(\bs{m}\cdot \overrightarrow{BD}=0\) 且 \(\bs{m}\cdot \overrightarrow{BF}=0\)，
即 \(-x+y=0\)，\(2y+hz=0\).

不妨令 \( y=1 \)，可得
\(\bs{m}=\left( 1,1,-\frac{2}{h} \right)\).

由题意，
\(\left| \cos \angle (\bs{m},\bs{n}) \right|
=\frac{|\bs{m}\cdot \bs{n}|}{|\bs{m}||\bs{n}|}
=\frac{\left| 4-\frac{2}{h} \right|}{3\sqrt{2+\frac{4}{h^2}}}
=\frac13\).

解得
\(h=\frac87\).

经检验，符合题意.

所以，线段 \( CF \) 的长为
\(\frac87\).
\end{solution}
\begin{problem}
设椭圆 \( \dfrac{x^2}{a^2}+\dfrac{y^2}{b^2}=1 \)（\( a>b>0 \)）的左焦点为 \( F \)，上顶点为 \( B \). 已知椭圆的短轴长为 \( 4 \)，离心率为 \( \dfrac{\sqrt5}{5} \).

\begin{enumerate}
\item 求椭圆的方程；
\item 设点 \( P \) 在椭圆上，且异于椭圆的上、下顶点，点 \( M \) 为直线 \( PB \) 与 \( x \) 轴的交点，点 \( N \) 在 \( y \) 轴的负半轴上. 若 \( |ON|=|OF| \)（\( O \) 为原点），且 \( OP\perp MN \)，求直线 \( PB \) 的斜率.
\end{enumerate}

\end{problem}

\begin{solution}
（1）设椭圆的半焦距为 \( c \)，依题意，
\(2b=4\)，\(\frac{c}{a}=\frac{\sqrt5}{5}\)，
又
\(a^2=b^2+c^2\).

可得
\(a=\sqrt5\)，\(b=2\)，\(c=1\).

所以，椭圆方程为
\(\frac{x^2}{5}+\frac{y^2}{4}=1\).

（2）由题意，设
\(P(x_P,y_P)\ (x_P\ne 0)\)，\(M(x_M,0)\).

设直线 \( PB \) 的斜率为 \( k \)（\( k\ne 0 \)），又 \( B(0,2) \)，则直线 \( PB \) 的方程为
\(y=kx+2\).

与椭圆方程联立，得
\(\left\{\begin{array}{l}
y=kx+2,\\
\dfrac{x^2}{5}+\dfrac{y^2}{4}=1
\end{array}\right.\)
整理得
\((4+5k^2)x^2+20kx=0\).
可得
\(x_P=-\frac{20k}{4+5k^2}\).

代入 \( y=kx+2 \) 得
\(y_P=\frac{8-10k^2}{4+5k^2}\).

进而直线 \( OP \) 的斜率
\(\frac{y_P}{x_P}=\frac{4-5k^2}{-10k}\).

在 \( y=kx+2 \) 中，令 \( y=0 \)，得
\(x_M=-\frac{2}{k}\).

由题意得 \( N(0,-1) \)，所以直线 \( MN \) 的斜率为
\(-\frac{k}{2}\).

由 \( OP\perp MN \)，得
\(\frac{4-5k^2}{-10k}\cdot \left( -\frac{k}{2} \right)=-1\).

化简得
\(k^2=\frac{24}{5}\).
从而
\(k=\pm \frac{2\sqrt{30}}{5}\).

所以，直线 \( PB \) 的斜率为
\(\frac{2\sqrt{30}}{5}\)或\(-\frac{2\sqrt{30}}{5}\).
\end{solution}
\begin{problem}
设 \( \{a_n\} \) 是等差数列，\( \{b_n\} \) 是等比数列. 已知 \( a_1=4 \)，\( b_1=6 \)，\( b_2=2a_2-2 \)，\( b_3=2a_3+4 \).

\begin{enumerate}
\item 求 \( \{a_n\} \) 和 \( \{b_n\} \) 的通项公式；
\item 设数列 \( \{c_n\} \) 满足 \( c_1=1 \)，\( c_n=\left\{\begin{array}{ll}
1, & 2^k<n<2^{k+1},\\
b_k, & n=2^k,
\end{array}\right. \) 其中 \( k\in \N^* \).
\begin{enumerate}
\item 求数列 \( \{a_{2^n}(c_{2^n}-1)\} \) 的通项公式；
\item 求 \( \sum_{i=1}^{2^n} a_ic_i \)（\( n\in \N^* \)）.
\end{enumerate}
\end{enumerate}

\end{problem}

\begin{solution}
（1）设等差数列 \( \{a_n\} \) 的公差为 \( d \)，等比数列 \( \{b_n\} \) 的公比为 \( q \).

依题意得
\(\left\{\begin{array}{l}
6q=2(4+d)-2=6+2d,\\
6q^2=2(4+2d)+4=12+4d.
\end{array}\right.\)

解得
\(d=3\)，\(q=2\).

故
\(a_n=4+(n-1)\times 3=3n+1\)，\(b_n=6\times 2^{n-1}=3\times 2^n\).

所以，\( \{a_n\} \) 的通项公式为
\(a_n=3n+1\)，
\( \{b_n\} \) 的通项公式为
\(b_n=3\times 2^n\).

（2）（i）因为 \( c_{2^n}=b_n \)，所以
\(a_{2^n}(c_{2^n}-1)=a_{2^n}(b_n-1)=(3\times 2^n+1)(3\times 2^n-1)=9\times 4^n-1\).

所以，数列 \( \{a_{2^n}(c_{2^n}-1)\} \) 的通项公式为
\(a_{2^n}(c_{2^n}-1)=9\times 4^n-1\).

（2）（ii）\(\sum_{i=1}^{2^n} a_ic_i=\sum_{i=1}^{2^n}\left[ a_i+a_i(c_i-1) \right]=\sum_{i=1}^{2^n} a_i+\sum_{i=1}^{2^n} a_i(c_i-1)\).

又因为只有当 \( i=2^k \)（\( 1\le k\le n \)）时，\( c_i-1\ne 0 \)，故
\(\sum_{i=1}^{2^n} a_i(c_i-1)=\sum_{k=1}^{n} a_{2^k}(c_{2^k}-1)=\sum_{k=1}^{n}(9\times 4^k-1)\).

于是
\examdisplayaligned{
\sum_{i=1}^{2^n} a_ic_i&=\sum_{i=1}^{2^n} (3i+1)+\sum_{k=1}^{n}(9\times 4^k-1)\\
&=\left( 3\cdot \frac{2^n(2^n+1)}{2}+2^n \right)+9\sum_{k=1}^{n}4^k-n\\
&=3\cdot 2^{2n-1}+5\cdot 2^{n-1}+9\cdot \frac{4(4^n-1)}{3}-n\\
&=27\times 2^{2n-1}+5\times 2^{n-1}-n-12.
}.
\end{solution}
\begin{problem}
设函数 \( f(x)=\e^x\cos x \)，\( g(x) \) 为 \( f(x) \) 的导函数.

\begin{enumerate}
\item 求 \( f(x) \) 的单调区间；
\item 当 \( x\in \left[ \dfrac{\pi}{4},\dfrac{\pi}{2} \right] \) 时，证明 \( f(x)+g(x)\left( \dfrac{\pi}{2}-x \right)\ge 0 \)；
\item 设 \( x_n \) 为函数 \( u(x)=f(x)-1 \) 在区间 \( \left( 2n\pi+\dfrac{\pi}{4},2n\pi+\dfrac{\pi}{2} \right) \) 内的零点，其中 \( n\in \N \)，证明
\(2n\pi+\dfrac{\pi}{2}-x_n<\dfrac{\e^{-2n\pi}}{\sin x_0-\cos x_0}\).
\end{enumerate}

\end{problem}

\begin{solution}
（1）由已知，
\(f'(x)=\e^x(\cos x-\sin x)\).

当
 \(x\in \left( 2k\pi+\frac{\pi}{4},2k\pi+\frac{5\pi}{4} \right)\)（\(k\in \Z\)）
时，有 \( \sin x>\cos x \)，得 \( f'(x)<0 \)，则 \( f(x) \) 单调递减；

当
 \(x\in \left( 2k\pi-\frac{3\pi}{4},2k\pi+\frac{\pi}{4} \right)\)（\(k\in \Z\)）
时，有 \( \sin x<\cos x \)，得 \( f'(x)>0 \)，则 \( f(x) \) 单调递增.

所以，\( f(x) \) 的单调递增区间为
\(\left( 2k\pi-\frac{3\pi}{4},2k\pi+\frac{\pi}{4} \right)\)（\(k\in \Z\)），
\( f(x) \) 的单调递减区间为
\(\left( 2k\pi+\frac{\pi}{4},2k\pi+\frac{5\pi}{4} \right)\)（\(k\in \Z\)）.

（2）记
\(h(x)=f(x)+g(x)\left( \frac{\pi}{2}-x \right)\).

依题意及（1）有
\(g(x)=\e^x(\cos x-\sin x)\)，
从而
\(g'(x)=-2\e^x\sin x\).

当
 \(x\in \left( \frac{\pi}{4},\frac{\pi}{2} \right)\)
时，\( g'(x)<0 \)，故
\(h'(x)=f'(x)+g'(x)\left( \frac{\pi}{2}-x \right)-g(x)
=g'(x)\left( \frac{\pi}{2}-x \right)<0\).

因此，\( h(x) \) 在区间 \( \left[ \dfrac{\pi}{4},\dfrac{\pi}{2} \right] \) 上单调递减，进而
\(h(x)\ge h\left( \frac{\pi}{2} \right)=f\left( \frac{\pi}{2} \right)=0\).

所以，当 \( x\in \left[ \dfrac{\pi}{4},\dfrac{\pi}{2} \right] \) 时，
\(f(x)+g(x)\left( \frac{\pi}{2}-x \right)\ge 0\).

（3）依题意，
\(u(x_n)=f(x_n)-1=0\)，
即
\(\e^{x_n}\cos x_n=1\).

记
\(y_n=x_n-2n\pi\)，
则
\(y_n\in \left( \frac{\pi}{4},\frac{\pi}{2} \right)\).

且
\(f(y_n)=\e^{y_n}\cos y_n
=\e^{x_n-2n\pi}\cos (x_n-2n\pi)
=\e^{-2n\pi}\).

由
\(f(y_n)=\e^{-2n\pi}<1=f(y_0)\)
及（1）得
\(y_n>y_0\).

由（2）知，当 \(x\in \left( \frac{\pi}{4},\frac{\pi}{2} \right)\)
时，\( g'(x)<0 \)，所以 \( g(x) \) 在 \( \left[ \dfrac{\pi}{4},\dfrac{\pi}{2} \right] \) 上为减函数，

因此
\(g(y_n)<g(y_0)<g\left( \frac{\pi}{4} \right)=0\).

又由（2）知
\(f(y_n)+g(y_n)\left( \frac{\pi}{2}-y_n \right)\ge 0\)，
故
\(\frac{\pi}{2}-y_n\le -\frac{f(y_n)}{g(y_n)}
=-\frac{\e^{-2n\pi}}{g(y_n)}
<-\frac{\e^{-2n\pi}}{g(y_0)}\).

而
\(g(y_0)=\e^{y_0}(\cos y_0-\sin y_0)
=-\e^{y_0}(\sin y_0-\cos y_0)\).

又因为 \( y_0=x_0 \)，所以
\(\frac{\pi}{2}-y_n<\frac{\e^{-2n\pi}}{\e^{y_0}(\sin y_0-\cos y_0)}
<\frac{\e^{-2n\pi}}{\sin x_0-\cos x_0}\).

从而
\(2n\pi+\frac{\pi}{2}-x_n<\frac{\e^{-2n\pi}}{\sin x_0-\cos x_0}\).
\end{solution}
